% Mathématiques 6e — cours V3
% Source LaTeX rendue côté Astro/KaTeX.

\section{Objectifs}

Dans ce chapitre, on apprend à :

\begin{itemize}
\item lire et écrire correctement de grands nombres entiers et des nombres décimaux ;
\item comprendre la valeur d’un chiffre selon sa position ;
\item décomposer et recomposer un nombre ;
\item passer d’une écriture décimale à une écriture fractionnaire décimale ;
\item comparer, ranger, encadrer et arrondir des nombres décimaux ;
\item placer et repérer des nombres sur une demi-droite graduée ;
\item calculer avec des nombres décimaux ;
\item utiliser un ordre de grandeur pour vérifier qu’un résultat est plausible ;
\item choisir l’opération adaptée dans un problème.
\end{itemize}

\section{1. Situation de départ — le même chiffre peut avoir des valeurs différentes}

Dans les nombres

\[
4\,205{,}37
\]

et

\[
0{,}043,
\]

le chiffre $4$ n’a pas la même valeur.

Dans le premier nombre, il représente des milliers :

\[
4\times1\,000=4\,000.
\]

Dans le second, il représente des centièmes :

\[
4\times\dfrac{1}{100}=0{,}04.
\]

La valeur d’un chiffre dépend donc de \textbf{son rang} dans l’écriture du nombre.

C’est le principe fondamental de la numération décimale de position.

\section{2. Notre numération est décimale et positionnelle}

Notre système utilise dix chiffres :

\[
0,\ 1,\ 2,\ 3,\ 4,\ 5,\ 6,\ 7,\ 8,\ 9.
\]

Il est \textbf{décimal} car les rangs sont organisés par puissances de dix.

Il est \textbf{positionnel} car la valeur d’un chiffre dépend de sa position.

Les principaux rangs utilisés en 6e sont :

| Rang | Valeur d’une unité du rang |
|---|---:|
| millier | $1\,000$ |
| centaine | $100$ |
| dizaine | $10$ |
| unité | $1$ |
| dixième | $\dfrac{1}{10}=0{,}1$ |
| centième | $\dfrac{1}{100}=0{,}01$ |
| millième | $\dfrac{1}{1\,000}=0{,}001$ |

\coursefigure{/images/mathematiques/6eme/nombres-entiers-decimaux-position.svg}{Schéma du tableau de numération décimale séparant partie entière et partie décimale.}{Pour le nombre $4\,205{,}37$, les chiffres occupent successivement les rangs des milliers, centaines, dizaines, unités, dixièmes et centièmes.}

\section{3. Décomposer un nombre entier ou décimal}

Décomposer un nombre consiste à faire apparaître la valeur de chaque chiffre.

\subsection{Exemple développé 1 — décomposer un nombre}

Considérons :

\[
4\,205{,}37.
\]

On peut écrire :

\[
4\,205{,}37
=
4\times1\,000
+
2\times100
+
0\times10
+
5\times1
+
3\times\dfrac{1}{10}
+
7\times\dfrac{1}{100}.
\]

Donc :

\[
4\,205{,}37
=
4\,000
+
200
+
5
+
0{,}3
+
0{,}07.
\]

On peut aussi utiliser des fractions décimales :

\[
4\,205{,}37
=
4\,205+\dfrac{3}{10}+\dfrac{7}{100}.
\]

Ou encore :

\[
4\,205{,}37
=
\dfrac{420\,537}{100}.
\]

Ces écritures représentent toutes \textbf{le même nombre}.

\section{4. Passer d’un rang au rang voisin}

Deux rangs voisins sont liés par un facteur $10$.

Par exemple :

\[
1\text{ dizaine}=10\text{ unités},
\]

\[
1\text{ unité}=10\text{ dixièmes},
\]

\[
1\text{ dixième}=10\text{ centièmes}.
\]

On peut donc écrire :

\[
1=10\times\dfrac{1}{10}
=100\times\dfrac{1}{100}
=1\,000\times\dfrac{1}{1\,000}.
\]

Inversement :

\[
\dfrac{1}{10}=0{,}1,
\qquad
\dfrac{1}{100}=0{,}01,
\qquad
\dfrac{1}{1\,000}=0{,}001.
\]

Ces relations sont essentielles pour comprendre les multiplications et divisions par $10$, $100$ et $1\,000$.

\section{5. Reconnaître un nombre décimal}

Un nombre décimal est un nombre qui peut s’écrire sous la forme :

\[
\dfrac{a}{10^n},
\]

où $a$ est un entier et $n$ un entier naturel.

Par exemple :

\[
3{,}47=\dfrac{347}{100}.
\]

Ainsi $3{,}47$ est un nombre décimal.

De même :

\[
12=\dfrac{12}{1}
\]

est aussi un nombre décimal.

Un entier est donc un cas particulier de nombre décimal.

\section{6. Zéros inutiles et égalité de nombres décimaux}

Ajouter des zéros à droite de la partie décimale ne change pas la valeur du nombre.

Ainsi :

\[
4{,}7=4{,}70=4{,}700.
\]

En effet :

\[
4{,}7
=
4+\dfrac{7}{10},
\]

tandis que

\[
4{,}70
=
4+\dfrac{70}{100}
=
4+\dfrac{7}{10}.
\]

Les deux écritures désignent donc le même nombre.

Attention : on ne peut pas ajouter un zéro n’importe où.

Par exemple :

\[
4{,}7\neq4{,}07.
\]

\section{7. Comparer deux nombres décimaux}

Pour comparer deux nombres décimaux, on compare les rangs de gauche à droite.

\subsection{Méthode}

Pour comparer $8{,}09$ et $8{,}9$ :

\begin{enumerate}
\item les parties entières sont égales ;
\item on compare les dixièmes ;
\item on peut écrire $8{,}9=8{,}90$.
\end{enumerate}

Ainsi :

\[
8{,}09<8{,}90.
\]

Donc :

\[
\boxed{8{,}09<8{,}9}.
\]

Il ne faut pas comparer les parties décimales comme si elles étaient des nombres entiers séparés.

\section{8. Ranger plusieurs nombres}

Ranger dans l’ordre croissant signifie aller du plus petit au plus grand.

Par exemple, pour :

\[
5{,}8,\qquad5{,}72,\qquad5{,}079,\qquad5{,}81,
\]

on peut écrire :

\[
5{,}8=5{,}800,
\]

\[
5{,}72=5{,}720,
\]

\[
5{,}079=5{,}079,
\]

\[
5{,}81=5{,}810.
\]

On obtient alors :

\[
\boxed{
5{,}079<5{,}72<5{,}8<5{,}81
}.
\]

\section{9. La demi-droite graduée}

Une demi-droite graduée permet de représenter les nombres et de visualiser leur ordre.

\coursefigure{/images/mathematiques/6eme/nombres-entiers-decimaux-droite-graduee.svg}{Schéma d'une demi-droite graduée permettant de placer, comparer et encadrer des nombres décimaux.}{Si les points $A$, $B$ et $C$ ont respectivement pour abscisses $5{,}2$, $5{,}44$ et $5{,}7$, leur position permet de lire directement $5{,}2<5{,}44<5{,}7$.}

Pour placer correctement un nombre, il faut toujours déterminer la valeur d’une graduation.

Si l’intervalle entre $5$ et $6$ est partagé en dix parties égales, une graduation vaut :

\[
\dfrac{6-5}{10}=0{,}1.
\]

Chaque graduation correspond donc à un dixième.

\section{10. Encadrer et intercaler}

Encadrer un nombre consiste à trouver un nombre plus petit et un nombre plus grand.

Par exemple :

\[
5{,}8<5{,}83<5{,}9.
\]

C’est un encadrement au dixième.

On peut aussi encadrer au centième :

\[
5{,}83<5{,}837<5{,}84.
\]

Intercaler consiste à trouver un nombre entre deux nombres donnés.

Entre $2{,}3$ et $2{,}4$, on peut par exemple choisir :

\[
2{,}35.
\]

On a bien :

\[
2{,}3<2{,}35<2{,}4.
\]

\section{11. Arrondir un nombre décimal}

Arrondir consiste à remplacer un nombre par une valeur proche choisie à un rang donné.

\subsection{Arrondi à l’unité}

Pour arrondir $17{,}6$ à l’unité, on compare les deux entiers voisins :

\[
17<17{,}6<18.
\]

Le nombre $17{,}6$ est plus proche de $18$.

Donc :

\[
17{,}6\approx18.
\]

\subsection{Arrondi au dixième}

Pour arrondir $5{,}837$ au dixième, on observe le chiffre des centièmes.

On obtient :

\[
5{,}837\approx5{,}8.
\]

\subsection{Arrondi au centième}

On observe cette fois le chiffre des millièmes :

\[
5{,}837\approx5{,}84.
\]

L’arrondi dépend donc du rang demandé.

\section{12. Additionner et soustraire des nombres décimaux}

Pour poser une addition ou une soustraction de nombres décimaux, il faut \textbf{aligner les virgules}.

\subsection{Exemple}

Calculons :

\[
18{,}7+4{,}56.
\]

On peut écrire :

\[
18{,}70+4{,}56=23{,}26.
\]

Donc :

\[
\boxed{18{,}7+4{,}56=23{,}26}.
\]

Pour la soustraction :

\[
18{,}7-4{,}56=14{,}14.
\]

Le zéro ajouté à droite de $18{,}7$ ne change pas le nombre, mais facilite le calcul posé.

\section{13. Multiplier ou diviser par $10$, $100$ ou $1\,000$}

Multiplier par $10$ signifie rendre chaque chiffre dix fois plus grand en valeur.

Par exemple :

\[
3{,}47\times10=34{,}7.
\]

De même :

\[
3{,}47\times100=347,
\]

et :

\[
3{,}47\times1\,000=3\,470.
\]

Inversement :

\[
347\div10=34{,}7,
\]

\[
347\div100=3{,}47,
\]

\[
347\div1\,000=0{,}347.
\]

Il vaut mieux raisonner sur les \textbf{rangs de numération} que mémoriser une règle du type « déplacer la virgule ».

\section{14. Multiplier par $0{,}1$, $0{,}01$ et $0{,}001$}

Comme :

\[
0{,}1=\dfrac{1}{10},
\]

on a :

\[
a\times0{,}1=a\div10.
\]

Par exemple :

\[
42{,}5\times0{,}1=4{,}25.
\]

De même :

\[
42{,}5\times0{,}01=0{,}425,
\]

et :

\[
42{,}5\times0{,}001=0{,}0425.
\]

Ces calculs utilisent directement la structure décimale des nombres.

\section{15. Comprendre la multiplication de deux nombres décimaux}

La multiplication de deux nombres décimaux ne doit pas être vue comme une simple « addition répétée ».

Elle peut représenter, par exemple, l’aire d’un rectangle.

Un rectangle de longueur $3{,}2\ \text{m}$ et de largeur $1{,}5\ \text{m}$ a pour aire :

\[
3{,}2\times1{,}5.
\]

On peut décomposer :

\[
1{,}5=1+0{,}5.
\]

Donc :

\[
3{,}2\times1{,}5
=
3{,}2\times1
+
3{,}2\times0{,}5.
\]

Ainsi :

\[
3{,}2\times1{,}5
=
3{,}2+1{,}6
=
4{,}8.
\]

Le produit vaut donc :

\[
\boxed{4{,}8}.
\]

\section{16. Exemple développé 2 — calculer et contrôler un produit}

Calculons :

\[
48{,}7\times19{,}8.
\]

Avant le calcul exact, cherchons un ordre de grandeur.

On peut utiliser :

\[
48{,}7\approx50
\]

et

\[
19{,}8\approx20.
\]

Ainsi :

\[
48{,}7\times19{,}8
\approx
50\times20
=
1\,000.
\]

Le résultat exact devra donc être proche de $1\,000$.

Le produit exact est :

\[
48{,}7\times19{,}8=964{,}26.
\]

Ce résultat est cohérent avec l’ordre de grandeur.

Un résultat comme :

\[
96{,}426
\]

serait immédiatement suspect.

\section{17. Diviser un nombre décimal par un entier}

En 6e, on rencontre notamment des divisions d’un nombre décimal par un entier non nul inférieur à $10$.

\subsection{Exemple développé 3 — partage}

On répartit $7{,}5\ \text{L}$ de jus dans $6$ bouteilles identiques.

Le volume d’une bouteille est :

\[
7{,}5\div6.
\]

On obtient :

\[
7{,}5\div6=1{,}25.
\]

Chaque bouteille contient donc :

\[
\boxed{1{,}25\ \text{L}}.
\]

Le quotient doit être interprété dans le contexte.

\section{18. Division euclidienne}

La division euclidienne concerne des nombres entiers.

Diviser $157$ par $12$, c’est chercher un quotient entier $q$ et un reste $r$ tels que :

\[
157=12\times q+r,
\]

avec :

\[
0\le r<12.
\]

On trouve :

\[
157=12\times13+1.
\]

Le quotient est donc :

\[
q=13,
\]

et le reste :

\[
r=1.
\]

Cette division répond par exemple à une question du type : combien de groupes complets peut-on former ?

\section{19. Ordre de grandeur}

Un ordre de grandeur est une valeur simple et proche du résultat exact.

Il sert à vérifier la cohérence d’un calcul.

\subsection{Addition}

Pour :

\[
197{,}8+403{,}2,
\]

on peut estimer :

\[
200+400=600.
\]

\subsection{Multiplication}

Pour :

\[
29{,}7\times4{,}1,
\]

on peut estimer :

\[
30\times4=120.
\]

\subsection{Division}

Pour :

\[
39{,}8\div8,
\]

on peut estimer :

\[
40\div8=5.
\]

L’ordre de grandeur ne remplace pas le résultat exact : il permet de le contrôler.

\section{20. Pourcentage : une écriture parmi d’autres}

Le pourcentage est une écriture liée à une fraction de dénominateur $100$.

Ainsi :

\[
25\,\%=\dfrac{25}{100}=0{,}25.
\]

De même :

\[
50\,\%=\dfrac{50}{100}=0{,}5,
\]

et :

\[
100\,\%=1.
\]

Cette écriture sera largement utilisée dans les autres chapitres.

\section{21. Méthode — résoudre un problème numérique}

Pour résoudre un problème, on peut suivre six étapes.

\subsection{Étape 1 — identifier la question}

Que cherche-t-on ?

\subsection{Étape 2 — repérer les données utiles}

Toutes les informations de l’énoncé ne sont pas forcément nécessaires.

\subsection{Étape 3 — choisir l’opération}

Il faut décider si la situation relève d’une addition, d’une soustraction, d’une multiplication, d’une division ou de plusieurs opérations successives.

\subsection{Étape 4 — prévoir un ordre de grandeur}

Avant le calcul exact, on estime le résultat.

\subsection{Étape 5 — effectuer le calcul}

Le calcul doit être posé ou organisé proprement.

\subsection{Étape 6 — interpréter}

On donne une phrase réponse avec l’unité lorsqu’elle est nécessaire.

\section{22. Exemple développé 4 — problème à plusieurs opérations}

Une association achète $18$ carnets à $2{,}75\ €$ l’unité et paie avec $60\ €$.

Le coût total est :

\[
18\times2{,}75.
\]

On obtient :

\[
18\times2{,}75=49{,}50.
\]

La monnaie rendue vaut :

\[
60-49{,}50=10{,}50.
\]

Donc l’association doit recevoir :

\[
\boxed{10{,}50\ €}.
\]

Vérification par ordre de grandeur :

\[
18\times3\approx54.
\]

Un coût de $49{,}50\ €$ est donc plausible.

\section{23. Automatismes à maîtriser}

Les automatismes utiles dans ce chapitre sont notamment :

\begin{itemize}
\item les tables de multiplication ;
\item les doubles et moitiés usuels ;
\item les compléments à $10$, $100$ et $1\,000$ ;
\item les relations entre unités, dixièmes, centièmes et millièmes ;
\item les multiplications et divisions par $10$, $100$ et $1\,000$ ;
\item les équivalences :
\end{itemize}

\[
\dfrac{1}{10}=0{,}1,
\]

\[
\dfrac{1}{100}=0{,}01,
\]

\[
\dfrac{1}{1\,000}=0{,}001.
\]

\section{24. Erreurs fréquentes}

\subsection{Comparer seulement le nombre de chiffres après la virgule}

On pourrait croire que :

\[
4{,}12>4{,}8
\]

parce que $12>8$.

C’est faux.

En écrivant :

\[
4{,}8=4{,}80,
\]

on voit que :

\[
\boxed{4{,}12<4{,}8}.
\]

\subsection{Oublier d’aligner les virgules}

Dans une addition ou une soustraction posée, on aligne les \textbf{rangs de numération}, donc les virgules.

\subsection{Ajouter des zéros au mauvais endroit}

On a :

\[
3{,}4=3{,}40,
\]

mais :

\[
3{,}4\neq3{,}04.
\]

\subsection{Confondre quotient décimal et division euclidienne}

Une division décimale peut produire un quotient décimal.

Une division euclidienne produit un quotient entier et un reste entier.

\subsection{Faire confiance à la calculatrice sans vérifier}

Même un calcul correctement saisi peut être mal interprété.

Un ordre de grandeur permet de détecter une erreur grossière.

\section{25. À retenir}

La numération décimale est fondée sur la valeur de position.

Deux rangs voisins sont liés par un facteur $10$.

Les relations essentielles sont :

\[
1=10\times\dfrac{1}{10},
\]

\[
1=100\times\dfrac{1}{100},
\]

\[
1=1\,000\times\dfrac{1}{1\,000}.
\]

Pour comparer des nombres décimaux, on compare les chiffres rang par rang.

Pour contrôler un calcul, on utilise un ordre de grandeur.

Enfin, dans un problème, le calcul n’est qu’une étape : il faut aussi choisir l’opération, vérifier la vraisemblance du résultat et l’interpréter.
